%! TEX root = thesis.tex

\chapter{Requirements and Specifications}
\label{sec:requirements}

\section{Introduction}

LLM training can involve thousands of GPUs running over the period of a couple of months. During this time, energy consumption and thereby CO2 consumption is immense, which is one of the arguments against LLMs in general. The CO2 consumption could possibly be reduced if the model training stops during times where the CO2 intensity of the energy grid is high, and resumes when the CO2 intensity is lower. This change of CO2 intensity mostly depends on the usage of renewable energy during this period in the specified area.

The simulation goal is to quantify the trade-off between reduced CO$_2$eq emissions and increased wall-clock training duration, compared to uninterrupted baseline training.

\section{Scoring of Policies}

\subsection{Target Function}

The policies are scored based on a composite metric that captures the trade-off between CO$_2$ savings and time overhead relative to baseline training.

Primary key performance indicators (KPIs):

\begin{align*}
  \text{CO2\_savings\_pct} &= \frac{CO2_{\text{baseline}} - CO2_{\text{policy}}}{CO2_{\text{baseline}}} \times 100, \\
  \text{time\_overhead\_pct} &= \frac{T_{\text{policy}} - T_{\text{baseline}}}{T_{\text{baseline}}} \times 100.
\end{align*}

Composite score used for ranking policies:

\[
  \text{score} = \frac{\text{CO2\_savings\_pct}}{\max(\text{time\_overhead\_pct}, \, \epsilon)} \quad \text{with } \epsilon > 0.
\]

Optional constraint for practical recommendations that only consider policies whose time overhead stays within a given budget:

\[
  \text{time\_overhead\_pct} \leq \text{overhead\_budget\_pct}.
\]

\subsection{Policy Variables}

\begin{itemize}
  \item $\theta_{\text{pause}}$ (gCO$_2$eq/kWh): Pause threshold (training pauses if grid intensity is above threshold).
  \item $\theta_{\text{hyst}}$ (gCO$_2$eq/kWh): Resume threshold (training resumes if grid intensity is below threshold (Hysteresis)).
  \item $t_{\text{start}}$ (datetime): Start time of training run.
  \item $\mathit{region}$ (categorical): Grid region used for carbon-intensity time series.
  \item $\mathit{pause\_granularity}$ (categorical): Allowed pause points (e.g., checkpoint, epoch).
\end{itemize}

\subsection{Constants and Parameters}

\subsubsection{Fixed constants per training-run profile}

\begin{itemize}
  \item $\mathit{model\_params}$ (count): Number of model parameters.
  \item $\mathit{dataset\_tokens}$ (count): Number of training tokens.
  \item $\mathit{gpu\_count}$ (count): Number of GPUs in cluster.
  \item $\mathit{gpu\_power\_train}$ (W/GPU): Average active GPU power during training.
  \item $\mathit{gpu\_power\_pause}$ (W/GPU): Average paused/idle GPU power.
  \item $\mathit{pue}$ (ratio): Data-center power usage effectiveness.
  \item $\mathit{checkpoint\_overhead\_time}$ (s/checkpoint): Time penalty per pause/resume checkpoint.
  \item $\mathit{checkpoint\_overhead\_energy}$ (Wh/checkpoint): Extra energy per pause/resume cycle.
\end{itemize}

\subsubsection{Scenario parameters (swept during study)}

\begin{itemize}
  \item $\mathit{threshold\_set}$ (list(gCO$_2$eq/kWh)): Candidate pause thresholds.
  \item $\mathit{hysteresis\_set}$ (list(gCO$_2$eq/kWh)): Candidate resume thresholds.
  \item $\mathit{regions\_set}$ (list(region)): Regions included in study.
  \item $\mathit{start\_times\_set}$ (list(datetime)): Candidate start dates/times.
  \item $\mathit{historical\_years}$ (list(year)): Historical periods used for replay.
  \item $\mathit{overhead\_budget\_pct}$ (\%): Maximum acceptable time overhead.
\end{itemize}

\subsubsection{Internal factors (time-dependent exogenous inputs)}

\begin{itemize}
  \item $\mathit{co2\_intensity}(t)$ (gCO$_2$eq/kWh): Grid carbon intensity time series.
  \item $\mathit{electricity\_price}(t)$ (EUR/kWh): Optional electricity price series for secondary analysis.
  \item $\mathit{data\_availability\_flag}(t)$ (binary): Indicates missing/invalid external data points.
\end{itemize}

\subsubsection{Internal state variables (simulation state)}

\begin{itemize}
  \item $\mathit{training\_progress}$ (\%): Fraction of total training work completed.
  \item $\mathit{is\_paused}$ (binary): Current pause/resume state.
  \item $\mathit{elapsed\_wall\_time}$ (h): Accumulated wall-clock duration.
  \item $\mathit{accumulated\_emissions}$ (gCO$_2$eq): Total emissions accumulated during run.
  \item $\mathit{pause\_count}$ (count): Number of pause/resume cycles.
\end{itemize}

\section{Objectives of the Simulation Study}

\subsection{Obligatory: Carbon Intensity \& Temporal Trade-offs}

The study shall \textbf{quantify the trade-off} between carbon intensity thresholds (gCO$_2$eq/kWh) and total training duration (latency), establishing a Pareto frontier for carbon-aware scheduling.

\subsection{Obligatory: Impact Assessment on SOTA Architectures}

The study shall \textbf{benchmark the absolute and relative carbon reduction potential} across state-of-the-art (SOTA) model architectures, evaluating how model scale and complexity influence the efficacy of carbon-aware training policies.

\subsection{Wishful: Spatiotemporal Optimization}

The study shall \textbf{identify the optimal temporal and geographical windows} for model training by analyzing the intersection of seasonal grid variability, regional renewable energy penetration, and local time-of-use carbon signals.

\subsection{Wishful: Geospatial Grid Analysis}

The study shall \textbf{conduct a comparative regional analysis} to determine the most sustainable geographical locations for training, accounting for carbon intensity.

\subsection{Optional: Computational Granularity \& Resumption Logic}

The study shall \textbf{evaluate the optimal checkpointing granularity} (e.g., batch-level vs. epoch-level) to determine the ideal breakpoint for pausing and resuming training, balancing carbon savings against the energy overhead of frequent state-saving and re-initialization.

\section{Scenarios}

\centering
\resizebox{\textwidth}{!}{%
  \begin{tabular}{l|c|c|c|c|c|c|c}
    \toprule
    Description & Model & Threshold & Hysteresis & Regions & Start times & Historical & Overhead \\
    \midrule
    \makecell{Original Deepseek \\ training} & Deepseek & 350, 400, 450 & 300, 350, 400 & CN & \makecell{2025-01-01 \\ 2025-02-18 \\ 2025-06-01} & \makecell{2022, 2023 \\ 2024, 2025} & 200 \\
    \hline
    \makecell{Deepseek training in \\ Sweden} & Deepseek & 20, 25 & 15, 20 & SE & \makecell{2025-01-01 \\ 2025-02-18 \\ 2025-06-01} & \makecell{2022, 2023 \\ 2024, 2025} & 200 \\
    \hline
    \makecell{Deepseek training in \\ Germany} & Deepseek & 300, 350 & 250, 300 & DE & \makecell{2025-01-01 \\ 2025-02-18 \\ 2025-06-01} & \makecell{2022, 2023 \\ 2024, 2025} & 200 \\
    \hline
    \makecell{Deepseek training with \\ multiple regions} & Deepseek & 300, 350 & 250, 300 & DE, IT, US & \makecell{2025-01-01 \\ 2025-06-01 \\ 2025-08-08} & \makecell{2022, 2023 \\ 2024, 2025} & 150 \\
    \hline
    \makecell{Kimi training in \\ China} & Kimi & 350, 400, 450 & 300, 350, 400 & CN & \makecell{2025-01-01 \\ 2025-02-18 \\ 2025-06-01} & \makecell{2022, 2023 \\ 2024, 2025} & 200 \\
    \hline
    \makecell{Kimi training in \\ Sweden} & Kimi & 20, 25 & 15, 20 & SE & \makecell{2025-01-01 \\ 2025-02-18 \\ 2025-06-01} & \makecell{2022, 2023 \\ 2024, 2025} & 200 \\
    \hline
    \makecell{Kimi training in \\ Germany} & Kimi & 300, 350 & 250, 300 & DE & \makecell{2025-01-01 \\ 2025-02-18 \\ 2025-06-01} & \makecell{2022, 2023 \\ 2024, 2025} & 200 \\
    \hline
    \makecell{Kimi training with \\ multiple regions} & Kimi & 300, 350 & 250, 300 & DE, IT, US & \makecell{2025-01-01 \\ 2025-06-01 \\ 2025-08-08} & \makecell{2022, 2023 \\ 2024, 2025} & 150 \\
    \bottomrule
  \end{tabular}%
}

\section{Resources}

\begin{itemize}
  \item "Electricity Maps" API for CO2 intensity data: \url{https://app.electricitymaps.com/map/live/fifteen_minutes}
  \item LLM Providers to find out about the training duration and energy consumption during training
\end{itemize}

\subsection{DeepSeek V3 Technical Report}

We base the DeepSeek scenario inputs on the training details reported by
\textcite{deepseek-ai_deepseek-v3_2024}.

\begin{itemize}
  \item Training Pipeline: Pretraining on 14.8T tokens (~2.788M GPU hours on 2048 H800 GPUs), Long Context Extension (32K then 128K), Post-Training with SFT and RL.
  \item Simulation inputs: $\mathit{model\_params}=671\mathrm{B}$, $\mathit{dataset\_tokens}=14.8\mathrm{T}$, $\mathit{gpu\_count}=2048$, GPU power and PUE to be looked up.
\end{itemize}

\subsection{Kimi K2 Technical Report}

We base the Kimi scenario inputs on the training details reported by
\textcite{kimi_team_kimi_2025}.

\begin{itemize}
  \item Training Pipeline: Pretraining on 15.5T tokens, post-training with a multi-stage agentic data synthesis pipeline and RL.
  \item Simulation inputs: $\mathit{model\_params}=1\mathrm{T}$ total ($32\mathrm{B}$ active), $\mathit{dataset\_tokens}=15.5\mathrm{T}$, GPU power and PUE to be looked up.
\end{itemize}

\subsection{Llama-3.1 405B}

\begin{itemize}
  \item Training Pipeline: Pretraining on 15T tokens (~30.84M GPU hours), fine-tuning details as available from modelcard.
  \item Simulation inputs: $\mathit{model\_params}=405\mathrm{B}$, $\mathit{dataset\_tokens}=15\mathrm{T}$, GPU TDP examples: 700W, other constants TBD.
\end{itemize}

% End of requirements chapter
